Para \(0<\delta<1\), sea \(y_{\delta,\varepsilon}\) la solución del problema
\[
-y''=(1-\delta)y+\sin(2x)+\varepsilon\sin x,
\qquad 0<x<\pi,\qquad y(0)=y(\pi)=0.
\]
Se consideran dos familias de datos: \(\varepsilon=\delta\) y \(\varepsilon=\delta^2\). En ambas, \(\delta\) tiende a cero por valores positivos.

Obtenga las soluciones y justifique su unicidad para \(0<\delta<1\). Compare sus límites en todo \([0,\pi]\) y el problema que resulta al poner directamente \(\delta=\varepsilon=0\). ¿Determina ese problema límite una única aproximación válida para ambas familias? Justifique la respuesta a partir de las soluciones, precisando el papel de los términos que desaparecen.
